\section{Relative generation from residual maximality}
For a finite indexed subgroup family \(H_i\leq G\), write \(H_J=\bigcap_{j\in J}H_j\), with \(H_\varnothing=G\) and \(B=H_I\).

\textbf{Theorem UB.Relative (generation from residual maximality).} Suppose all \(H_J\), \(J\subseteq I\), are distinct, and for every \(J\subsetneq I\) there is \(k\notin J\) such that \(H_{J\cup\{k\}}\) is maximal in \(H_J\). Choose \begin{equation}
s_i\in H_{I\setminus\{i\}}\setminus H_i.
\label{eq:02-relative:6-2}
\end{equation} Then \begin{equation}
H_J=\langle B,s_i:i\notin J\rangle\quad(J\subseteq I).
\label{eq:02-relative:6-3}
\end{equation} In particular these elements generate \(G\) irredundantly relative to \(B\), and their relative parabolics satisfy \begin{equation}
\langle B,s_i:i\in S\rangle\cap\langle B,s_i:i\in T\rangle
=\langle B,s_i:i\in S\cap T\rangle.
\label{eq:02-relative:6-4}
\end{equation} The group need not be finite.

\emph{Proof.} Distinctness of \(H_{I\setminus\{i\}}\) and \(B\) makes \eqref{eq:02-relative:6-2} possible, since its intersection with \(H_i\) is \(B\). Every \(s_j\) with \(j\ne i\), and every element of \(B\), belongs to \(H_i\), while \(s_i\) does not. This proves relative independence.

Prove \eqref{eq:02-relative:6-3} by downward induction from \(J=I\). Choose \(k\) as in the hypothesis. By induction, \(H_{J\cup\{k\}}\) is generated by \(B\) and the \(s_i\) with \(i\notin J\cup\{k\}\). The element \(s_k\) belongs to \(H_J\), because \(J\subseteq I\setminus\{k\}\), and lies outside \(H_{J\cup\{k\}}\), because it lies outside \(H_k\). Maximality therefore says that adjoining \(s_k\) generates \(H_J\). This is \eqref{eq:02-relative:6-3}. Finally substitute \(J=I\setminus S\) and \(I\setminus T\) into the elementary equality \(H_J\cap H_{J'}=H_{J\cup J'}\) to obtain \eqref{eq:02-relative:6-4}. \(\square\)

\noindent\textbf{Theorem UB.Generation (the geometric interface).}
For a finite-rank firm geometry with a residually weakly primitive action,
choose one chamber and its type stabilizers $H_i$. Every choice
\[
 s_i\in H_{I\setminus\{i\}}\setminus H_i
\]
is independent relative to $B=H_I$, generates $G$ over $B$, and satisfies
\[
 H_J=\gen{B,s_i:i\notin J}\qquad(J\subseteq I).
\]
Every such geometry is residually connected. For finite $G$ its rank is
at most $\mu(G,B)$. Without residual weak primitivity, a firm
flag-transitive geometry still has a relatively independent family of
panel choices and, for finite $G$, rank at most $\mu'(G,B)$.
\begin{proof}
Lemma UB.Flags supplies flag transitivity and coset coordinates for the
first assertion. Firmness gives distinct indexed intersections by the
panel argument in Proposition UB.Coset. Residual weak primitivity is
exactly (G4), so Theorem UB.Relative applies to the displayed choices.
Its formula implies (G2): if at least two types remain, each of the
remaining generators belongs to an intersection obtained by fixing
another remaining type. Consequently Proposition UB.Coset gives residual
connectedness. The relative generating family has size equal to the
rank, proving the bound by $\mu(G,B)$.

For the last assertion no maximality is needed. Firmness and flag
transitivity again give panel choices. For a fixed $i$, all other choices
and $B$ lie in $H_i$, whereas $s_i$ does not. They are therefore relatively
independent, which gives the bound by $\mu'(G,B)$.
\end{proof}
